\section{Primitive certification and operational recurrence exhaustion}\label{sec:algebraic}
The exact criterion in Theorem~\ref{thm:transfer} does not require finite-dimensional physical dynamics. An effective consequence does. In this section all parameter and program sets are compact semialgebraic sets defined over the real algebraic numbers, all alphabets are finite, and primitive marked-response matrix entries and task readouts are continuous semialgebraic functions. Thus the finite-response continuity hypotheses remain in force. They may be rational after introducing algebraic auxiliary variables. A uniform primitive termination certificate is part of the input. This is not an algorithm for arbitrary computable Borel kernels.

\begin{lemma}[Finite instrumental resolvents]\label{lem:resolvent}
For a finite classical physical experiment with retained labels, let $A$ be the substochastic alive matrix on physical-state--label pairs, $B$ the one-step exit matrix to boundary labels, and $E$ the preparation matrix. Assume $\rho(A)<1$ uniformly. If $v$ is the expected one-step score from an alive pair, then
\begin{equation}\label{eq:resolvents}
 P=E(I-A)^{-1}B,\qquad
 T=E(I-A)^{-2}B,\qquad
 r=E(I-A)^{-1}v.
\end{equation}
The same formulas hold for finite-dimensional quantum instruments after replacing $A$ by the completely positive alive superoperator on the direct sum of Hermitian physical blocks indexed by observer labels, and replacing matrix rows by preparation states and exit/score functionals. In either case the excursion arrays are semialgebraic functions of the primitive data.
\end{lemma}
\begin{proof}
An excursion exiting after $j\ge1$ calls contributes $EA^{j-1}B$ to $P$ and $jEA^{j-1}B$ to $T$. The absolutely convergent sums are $(I-A)^{-1}$ and $(I-A)^{-2}$; the reward sum is the alive resolvent applied to the one-step score. In the quantum case work with unnormalized states. Composition and summation of marked completely positive maps are linear on the finite real vector space of Hermitian blocks. An exit functional applied to $A^{j-1}$ has exactly the same duration weight $j$. A uniform trace-decay certificate gives absolute convergence. Inversion of a finite nonsingular real matrix is rational in its entries. No normalization at a zero-probability report or closure of an unbounded operator is involved.
\end{proof}
A fixed preparation call can be included as an extra transient block; it must not be removed from $T$ or $r$. A positive classical or quantum Lyapunov functional giving $\norm{A^j}\le C\alpha^j$, $\alpha<1$, is a sufficient uniform termination input and supplies every finite return moment.

\begin{theorem}[Decidable task regularity and algebraic modulus]\label{thm:algebraic}
For the finite algebraic instruments just described, gain continuity on a specified compact semialgebraic family $X$ is decidable. If it holds, there exist effectively obtainable constants $a_0>0$, $C<\infty$ and rational $\kappa\ge0$ such that
\begin{equation}\label{eq:algebraic-price}
 \calC_X(a)\le C a^{-\kappa},\qquad 0<a<a_0.
\end{equation}
Consequently the acquisition error is
\begin{equation}\label{eq:algebraic-rate}
 O\bigl(N^{-1/(\kappa+1)}\bigr)+B_N,
\end{equation}
and the discount error is $O((1-\beta)^{1/(\kappa+1)})$ plus its tail term. The optimal robust average value and whether it is attained on the declared semialgebraic program class are decidable even when gain continuity fails. These conclusions are decision and global certification results, not polynomial-time synthesis claims.
\end{theorem}
\begin{proof}
By Lemma~\ref{lem:resolvent}, $P,T,r,S,c$ are semialgebraic. The ergodic projection is semialgebraic on each rank stratum. One direct formula is the coefficient of $z^{-1}$ in $(zI+I-S)^{-1}$: if the determinant has first nonzero term $a_qz^q$, semisimplicity at zero of $I-S$ makes that coefficient the coefficient of $z^{q-1}$ in the adjugate divided by $a_q$. There are only finitely many possible $q$. Thus the graph of $g=\Pi_Sc$ is a finite semialgebraic union. Continuity can be expressed by the first-order condition
\[
 \forall x\in X\ \forall\epsilon>0\ \exists\eta>0\
 \forall y\in X:\quad \norm{x-y}<\eta\Longrightarrow
 \norm{g(x)-g(y)}<\epsilon.
\]
Quantifier elimination over real closed fields decides it.

For the price, impose $u_1=0$, introduce upper and lower bounds for its coordinates, and encode the residual inequality in~\eqref{eq:corrector-price}. The infimum of the resulting finite linear-inequality problem and its supremum over $X$ have semialgebraic graphs, by quantifier elimination; an infimum need not be selected to make this statement. Under continuity they are finite for every $a>0$ by Theorem~\ref{thm:transfer}. A finite one-variable semialgebraic function has polynomial growth at an endpoint. The semialgebraic growth theorem, together with effective real-algebraic decomposition, supplies~\eqref{eq:algebraic-price} with rational exponent and algebraic bounds~\cite{BasuPollackRoy}. Balancing $a$ with $Ca^{-\kappa}/N$ proves~\eqref{eq:algebraic-rate}; the same balance works for discount.

Finally, the graph of the worst-model gain is semialgebraic using quantified upper bounds and arbitrarily close lower witnesses. Its infimum on the program set is therefore a real algebraic number when all defining coefficients are algebraic. The sentence asserting a program with that value decides attainment. The argument applies to a discontinuous objective as well, but then asserts neither compact attainment nor uniform finite-acquisition convergence unless the corresponding criterion holds. The algebraic assertions use standard real-algebraic machinery, not a numerical sample of the parameter continuum.
\end{proof}
The exponent in~\eqref{eq:algebraic-price} is a certified possible exponent, not necessarily the sharp exponent. The example in Section~\ref{sec:sharpness} identifies the exact price order directly. Algebraic descriptions of row probabilities are finite programs for real-arithmetic sampling; conversion to a particular bit-level random sampler, its variable running time and its approximation errors are separate implementation obligations.

\subsection{An actionable exhaustion beyond compact attainment}
A different conclusion remains available when the average objective is discontinuous. It must not be confused with Theorem~\ref{thm:algebraic}'s compatible case.

\begin{theorem}[Polynomial near-optimal recurrence schedule]\label{thm:schedule}
Assume the finite algebraic setting above, a uniform bound $\bar\mu$ on mean physical duration and a uniform return envelope over all legal programs. Suppose also that for each fixed program $\gamma$, the boundary support graph of $P(\theta,\gamma)$ is independent of $\theta\in\Theta$. Let
\[
 V=\inf_{\gamma\in\Gamma_M}\sup_{\theta\in\Theta}g(\theta,\gamma).
\]
There are $a_0>0$, $C<\infty$, rational $\kappa\ge0$ and a semialgebraic choice of common programs $\gamma_a$, $0<a<a_0$, such that
\begin{equation}\label{eq:schedule}
 \sup_\theta g(\theta,\gamma_a)\le V+a,
 \qquad
 \sup_\theta\norm{G_{P(\theta,\gamma_a)}}_\infty\le Ca^{-\kappa}.
\end{equation}
In particular the same program $\gamma_a$ satisfies, for every $N$,
\begin{equation}\label{eq:schedule-physical}
 \sup_\theta J_N(\theta,\gamma_a)
 \le V+a+2\bar\mu Ca^{-\kappa}/N+B_N,
\end{equation}
with the analogous discount bound. The data determine whether the unrestricted average infimum is attained. Without task compatibility, no assertion is made that the unrestricted finite-$N$ optimal values converge to $V$.
\end{theorem}
\begin{proof}
For fixed $\gamma$, the common graph fixes the number of closed classes as $\theta$ varies. The spectral-contour argument in Corollary~\ref{cor:group}, applied on compact $\Theta$, makes
\[
 H(\gamma)=\max_\theta\norm{G_{P(\theta,\gamma)}}_\infty
\]
finite. Across programs this function can be unbounded and discontinuous, but it is semialgebraic by the finite rank-stratum formulas and quantifier elimination. The set of pairs $(a,\gamma)$ satisfying the first inequality of~\eqref{eq:schedule} is nonempty for every $a>0$ and semialgebraic. Semialgebraic choice supplies a curve $\gamma_a$. The finite semialgebraic function $H(\gamma_a)$ has polynomial growth as $a\downarrow0$, proving the second inequality after shrinking $a_0$. Corollary~\ref{cor:group} gives~\eqref{eq:schedule-physical}. Attainment is the first-order statement proved decidable in Theorem~\ref{thm:algebraic}.
\end{proof}
This theorem converts recurrence exhaustion into a primitive-data, near-optimal implementation schedule. It permits the required recurrence level to diverge and does not manufacture an optimizer at the limit. The common-support assumption is testable in the finite algebraic setting; it is not a consequence of continuity alone. In particular, the rank-changing rare-herald family below uses Theorem~\ref{thm:transfer}, not this common-support theorem. Semialgebraic strategy asymptotics in stochastic games have classical precedents~\cite{BewleyKohlberg,BewleyKohlbergStationary}; the schedule here applies them to one fixed-memory robust causal program and its separately charged physical-return bridge.
